\setcounter{ex}{0}
\setcounter{paragraph}{0}
\PhanI
\begin{ex} 
(QG - 2022) Môt con lắc đơn chiều dài $\ell $ đang dao động điều hòa với biên độ góc $\alpha_0\ (rad) $. Biên độ dao động của con lắc là
\choice
{\True $s_0=\ell \alpha_0$}
{$s_0=\dfrac{\ell}{\alpha_0}$}
{$s_0=\ell^2\alpha_0$}
{$s_0=\dfrac{\alpha_0}{\ell}$}
\loigiai{
Biên độ dao động cung của con lắc là $s_0=\ell \alpha_0$}
\end{ex} \haidong{20}
%\loai{Công thức}
\begin{ex}
(QG - 2021) Một con lắc đơn đang dao động điều hòa với phương trình $s = S_0\cos(\omega t + \varphi)\ (S_0> 0)$. Đại lượng $S_0$ được gọi là
\choice
{\True biên độ của dao động}
{tần số của dao động}
{li độ góc của dao động}
{pha ban đầu của dao động}
\end{ex} \haidong{20}




%\loai{Dao động nhỏ con lắc đơn}
\begin{ex} %[3P1YF-1][Loai1]  
Dao động của con lắc đơn được xem là dao động điều hoà khi
\choice
{\True không có ma sát và dao động với biên độ nhỏ}
{biên độ dao động nhỏ}
{chu kì dao động không đổi}
{không có ma sát}
\loigiai{
Dao động của con lắc đơn được xem là dao động điều hòa khi không có ma sát và biên độ dao động là nhỏ.
}
\end{ex} \haidong{20}
\begin{ex} %[3P1YF-1][Loai1]  
Một con lắc đơn được treo tại một điểm cố định. Kéo con lắc ra khỏi vị trí cân bằng để dây treo hợp với phương thẳng đứng góc $60^\circ$ rồi buông nhẹ. Bỏ qua mọi ma sát và lực cản. Chuyển động của con lắc là
\choice
{chuyển động thẳng đều}
{\True dao động tuần hoàn}
{chuyển động tròn đều}
{dao động điều hòa}
\loigiai{
Dao động của con lắc đơn có biên độ góc lớn hơn $10^\circ$ là dao động tuần hoàn.
}
\end{ex} \haidong{20}

%\loai{Các yếu tố chu kì phụ thuộc}

\begin{ex} %[3P1BF-1][Loai1]  
Khi khối lượng của vật nặng là m thì chu kì dao động của con lắc đơn là T. Vậy khi tăng khối lượng của vật nặng lên 4 lần thì chu kì dao động của con lắc đơn lúc này là
\choice
{2T}
{\True T}
{4T}
{$\dfrac{T}{2}$}
\loigiai{
Chu kì của con lắc đơn được tính theo công thức: $T=2\pi \sqrt{{\dfrac{\ell}{g}}}$ không phụ thuộc vào khối lượng của vật nặng nên chu kì của con lắc đơn sau khi thay đổi khối lượng vật nặng vẫn là T.
}
\end{ex} \haidong{20}

\begin{ex}
Một con lắc đơn có chiều dài dây treo không đổi, thực hiện những dao động điều hoà ở một nơi có gia tốc trọng trường không đổi. Khi khối lượng vật nặng là $m_1$, chu kì con lắc là $T_1$. Khi khối lượng vật nặng là $m_2$, chu kì con lắc là $T_2$. Cả hai trường hợp con lắc đều dao động điều hoà. Hệ thức nào sau đây đúng?
\choice
{$\dfrac{T_1}{T_2}=\dfrac{m_1}{m_2}$}
{$\dfrac{T_1}{T_2}=\sqrt{\dfrac{m_1}{m_2}}$}
{$\dfrac{T_1}{T_2}=\dfrac{m_2}{m_1}$}
{\True $T_1=T_2$}
\end{ex} \haidong{20}
%\loai{Các đại lượng tỉ lệ}
\begin{ex} %[3P1YF-1][Loai1]  
Tại một nơi xác định, chu kì dao động điều hòa của con lắc đơn tỉ lệ thuận với
\choice
{gia tốc trọng trường}
{chiều dài con lắc}
{\True căn bậc hai chiều dài con lắc}
{căn bậc hai gia tốc trọng trường}
\loigiai{
Tại một nơi xác định, chu kì dao động của con lắc đơn tỉ lệ thuận với căn bậc hai chiều dài con lắc.
}
\end{ex} \haidong{20}
%\loai{Thay đổi thông số con lắc}
\begin{ex} %[3P1BF-1][Loai1]  
Chu kì dao động điều hòa của con lắc đơn sẽ giảm khi
\choice
{tăng chiều dài dây treo}
{giảm khối lượng vật nhỏ}
{giảm biên độ dao động}
{\True gia tốc trọng trường tăng}
\loigiai{
Chu kì dao động của con lắc đơn $T=2\pi \sqrt{\dfrac{\ell }{g}} \Rightarrow $ T giảm khi g tăng.
}
\end{ex} \haidong{20}

\begin{ex} %[3P1BF-1][Loai1]  
Để giảm tần số dao động con lắc đơn 2 lần, ta cần
\choice
{giảm chiều dài của dây 2 lần}
{giảm chiều dài của dây 4 lần}
{tăng chiều dài của dây 2 lần}
{\True tăng chiều dài của dây 4 lần}
\loigiai{
\begin{itemize}
\item Ta có: $f = \dfrac{1}{2\pi} \sqrt{\dfrac{g}{\ell}}$.
\item Ta thấy, để giảm tần số 2 lần thì cần tăng l lên 4 lần hoặc giảm g đi 4 lần.
\end{itemize}
}

\end{ex} \haidong{20}
%\loai{So sánh dao động hai con lắc}

\begin{ex} %[3P1BF-1][Loai1]  
Hai con lắc đơn có chiều dài lần lượt $\ell_1$ và $\ell_2$ với $\ell_1$ = 2$\ell_2$ dao động tự do tại cùng một vị trí trên trái đất, hãy so sánh tần số dao động của hai con lắc
\choice
{$f_1 = 2f_2$}
{$f_1 = \dfrac{f_2}{2}$}
{\True $f_2 = \sqrt{2}f_1$}
{$f_1 = \sqrt{2}f_2$}
\loigiai{
Ta có: $\dfrac{f_2}{f_1} = \sqrt{\dfrac{\ell_1}{\ell_2}} = \sqrt{2}\Rightarrow f_2 = \sqrt{2}.f_1$
}
\end{ex} \haidong{20}















%\loai{Lí thuyết vận tốc}

%\loai{Lí thuyết năng lượng}
\begin{ex} %[3P1-16.1B][Loai1]   
Với gốc thế năng được chọn tại vị trí cân bằng. Phát biểu nào sau đây là \textbf{sai} khi nói về cơ năng của con lắc đơn khi dao động điều hòa?
\choice
{Cơ năng bằng thế năng của vật ở vị trí biên}
{Cơ năng bằng tổng động năng và thế năng của vật khi qua vị trí bất kì}
{\True Cơ năng của con lắc đơn tỉ lệ thuận với biên độ góc}
{Cơ năng bằng động năng của vật khi qua vị trí cân bằng}
\loigiai{
Cơ năng của con lắc đơn không tỉ lệ thuận với biên độ góc $\Rightarrow $ C sai.
}
\end{ex} \haidong{20}
\begin{ex} %[3P1-16.1B][Loai1]  

Một con lắc đơn dao động điều hòa. Năng lượng sẽ thay đổi như thế nào nếu độ cao cực đại của vật tính từ vị trí cân bằng tăng 2 lần
\choice
{\True Tăng 2 lần}
{Giảm 2 lần}
{Tăng 4 lần}
{Giảm 4 lần}
\loigiai{
Ta có: $W = W_{t\max} = mgh\Rightarrow$ h tăng 2 lần thì W tăng 2 lần.
}
\end{ex} \haidong{20}


\begin{ex} %[3P1-16.1B][Loai1]  
Thế năng của con lắc đơn dao động điều hoà
\choice
{\True bằng với năng lượng dao động khi vật nặng ở biên}
{cực đại khi vật qua vị trí cân bằng}
{luôn không đổi vì quỹ đạo của vật được coi là thẳng}
{không phụ thuộc góc lệch của dây treo}
\loigiai{
\begin{itemize}
\item Ta có: $W_t = mg\ell(1-\cos\alpha)= \dfrac{mg\ell\alpha^2}{2}$, với $\alpha$ là góc lệch nhỏ
\item Thế năng của con lắc đơn dao động điều hòa phụ thuộc vào góc lệch của dây treo nên D, C sai.
\end{itemize}
}
\end{ex} \haidong{20}
\begin{ex} %[3P1-16.1Y][Loai1]   
Một con lắc đơn gồm dây treo dài $\ell$ và vật có khối lượng là m. Con lắc treo tại nơi có gia tốc rơi tự do là g. Kích thích con lắc dao động điều hòa với biên độ góc $\alpha _0.$ Biểu thức năng lượng dao động của con lắc là
\choice
{$2mg\ell\alpha _0^2$}
{\True $\dfrac{1}{2}{mg\ell}\alpha _0^2$}
{$mg\ell\alpha _0^2$}
{$\dfrac{2mg}{\ell}\alpha _0^2$}
\loigiai{
Cơ năng của con lắc được xác định bằng biểu thức: $W=\dfrac{1}{2}{mg\ell}\alpha _0^2$.
}
\end{ex} \haidong{20}
\begin{ex} %[3P1-16.1Y][Loai1]   
Một con lắc đơn có dây treo dài $\ell$, vật nặng khối lượng m đặt tại nơi có gia tốc trọng trường g, biên độ góc là $\alpha _0\leq 10^\circ$. Bỏ qua mọi lực cản. Chọn mốc thế năng là vị trí thấp nhất của vật. Khi con lắc đi qua vị trí có li độ góc $\alpha $ thì thế năng của vật nặng là
\choice
{$W_t=\dfrac{1}{2}{mg}\ell \alpha $}
{$W_t=\dfrac{1}{2}{g}\ell \alpha $}
{$W_t=\dfrac{1}{2}{g}\ell \alpha ^2$}
{\True $W_t=\dfrac{1}{2}{mg}\ell \alpha ^2$}
\loigiai{
Thế năng của vật nặng tại vị trí có li độ góc $\alpha $ được xác định bởi biểu thức $W_t=0,5mg\ell\alpha ^2$.
}
\end{ex} \haidong{20}
\begin{ex} %[3P1-16.1B][Loai1]   
(QG - 2011) Một con lắc đơn dao động điều hòa với biên độ góc $\alpha_0$. Lấy mốc thế năng ở vị trí cân bằng. Ở vị trí con lắc có động năng bằng thế năng thì li độ góc của nó bằng
\choice
{$\pm \dfrac{\alpha _0}{\sqrt{3}}$}
{$\pm \dfrac{\alpha _0}{2}$}
{\True $\pm \dfrac{\alpha _0}{\sqrt{2}}$}
{$\pm \dfrac{\alpha _0}{3}$}
\loigiai{Ở vị trí con lắc có động năng bằng thế năng thì li độ góc của nó bằng:
\begin{align*}
W_t=W_{\text{đ}}=\dfrac{W}{2}\Leftrightarrow \dfrac{mg\alpha^2}{2}=\dfrac{1}{2}\dfrac{mg\alpha_0^2}{2}\Rightarrow \alpha=\pm \dfrac{\alpha_0}{\sqrt{2}}
\end{align*}}
\end{ex} \haidong{20}

\begin{ex} %[3P1-16.1B][Loai1]   
(ĐH - 2010) Tại nơi có gia tốc trọng trường g, một con lắc đơn dao động điều hòa với biên độ góc $\alpha _0$ nhỏ. Lấy mốc thế năng ở vị trí cân bằng. Khi con lắc chuyển động nhanh dần theo chiều dương đến vị trí có động năng bằng thế năng thì li độ góc $\alpha $ của con lắc bằng
\choice
{$-\dfrac{\alpha _0}{\sqrt{3}}$}
{$\dfrac{\alpha _0}{\sqrt{2}}$}
{\True $-\dfrac{\alpha _0}{\sqrt{2}}$}
{$\dfrac{\alpha _0}{\sqrt{3}}$}
\loigiai{Con lắc đang đi theo chiều dương về trị trí cân bằng $\Rightarrow \alpha <0$
\begin{align*}
W_t=W_{\text{đ}}=\dfrac{W}{2}\Rightarrow \alpha =\pm \dfrac{\alpha _0}{\sqrt{2}} \Rightarrow \alpha =-\dfrac{\alpha _0}{\sqrt{2}
}.
\end{align*}}
\end{ex} \haidong{20}


\begin{ex} %[3P1-16.1B][Loai1]  

Tại nơi có gia tốc trọng trường g, một con lắc đơn dao động điều hòa với biên độ góc $\alpha_0$ nhỏ. Lấy mốc thế năng ở vị trí cân bằng. Khi con lắc chuyển động chậm dần theo chiều âm đến vị trí có động năng bằng ba lần thế năng thì li độ góc $\alpha$ của con lắc bằng
\choice
{$\dfrac{-\alpha_0}{\sqrt{3}}$}
{\True $\dfrac{-\alpha_0}{2}$}
{$\dfrac{\alpha_0}{2}$}
{$\dfrac{\alpha_0\sqrt{3}}{2}$}
\loigiai{
\begin{itemize}
\item Khi $W_{\text{đ}} = 3W_t \Rightarrow W_t = \dfrac{1}{4}W \Rightarrow \dfrac{1}{2}m\omega^2x^2 = \dfrac{1}{4}\dfrac{1}{2}m\omega^2A^2 \Rightarrow x = \pm \dfrac{A}{2} \Rightarrow \alpha = \pm \dfrac{\alpha_0}{2}$
\item Để vật chuyển động chậm dần theo chiều âm thì $\alpha = - \dfrac{\alpha_0}{2}$.
\end{itemize}
}
\end{ex} \haidong{20}
\setcounter{ex}{0}
\PhanII
\begin{ex}
Một con lắc đơn được gắn vật nặng $m$, chiều dài dây $\ell$ được kích thích dao động với góc bé hơn $10^\circ$.
\choiceTF[t]
{\True Có thể coi dao động đó là dao động điều hòa}
{Chu kì dao động phụ thuộc vào khối lượng $m$ của vật nặng}
{\True Khi tăng chiều dài dây thì chu kì dao động vật tăng}
{Chu kì dao động không phụ thuộc vào gia tốc trọng trường}
\loigiai{
\begin{itemchoice}
\itemch
\itemch
\itemch
\itemch
\end{itemchoice}
}
\end{ex} \haidong{20}
\begin{ex}
Đối với một con lắc đơn dao động nhỏ (biên độ góc không vượt quá $10^\circ$).
\choiceTF[t]
{\True Chu kì dao động tỉ lệ với căn bậc hai của chiều dài dây treo}
{Khi tăng khối lượng vật nặng lên 2 lần, chu kì tăng 2 lần}
{\True Con lắc chỉ dao động điều hoà khi bỏ qua mọi ma sát và biên độ đủ nhỏ}
{Chu kì dao động nhỏ hoàn toàn không phụ thuộc vào gia tốc trọng trường}
\loigiai{
\begin{itemchoice}
\itemch Chu kì $T = 2\pi\sqrt{\dfrac{\ell}{g}}\ \Rightarrow$ Đúng
\itemch Công thức $T$ không chứa khối lượng $m\ \Rightarrow$ Sai
\itemch Điều kiện dao động điều hoà của con lắc đơn là biên độ nhỏ và không ma sát\ \Rightarrow\ Đúng
\itemch $T$ phụ thuộc $g$ nên khẳng định không phụ thuộc $g$ là Sai
\end{itemchoice}
}
\end{ex}
\begin{ex}
Một con lắc đơn dao động điều hoà nhỏ, lấy mốc thế năng tại vị trí cân bằng.
\choiceTF[t]
{\True Ở vị trí biên, vận tốc bằng $0$ và thế năng đạt cực đại}
{Tổng động năng và thế năng biến thiên tuần hoàn theo thời gian}
{\True Cơ năng của con lắc phụ thuộc vào khối lượng quả nặng và biên độ góc}
{Khi vật nặng qua vị trí cân bằng, thế năng đạt giá trị lớn nhất}
\loigiai{
\begin{itemchoice}
\itemch Tốc độ $v = 0$ tại biên, $W_t$ cực đại.
\itemch Cơ năng $W$ không đổi theo thời gian.
\itemch $W = \dfrac{1}{2}mg\ell\alpha_0^{2}$ phụ thuộc $m,\ \alpha_0$.
\itemch Tại cân bằng $x = 0$ nên $W_t = 0$.
\end{itemchoice}
}
\end{ex}
\begin{ex}
Một con lắc đơn dao động nhỏ tại cùng một địa điểm trên Trái Đất.
\choiceTF[t]
{\True Khi chiều dài dây tăng 4 lần, chu kì dao động tăng 2 lần}
{Chu kì dao động nhỏ luôn phụ thuộc vào biên độ góc, dù biên độ rất nhỏ}
{\True Hai con lắc đơn có cùng chiều dài sẽ dao động với cùng chu kì}
{Đưa con lắc lên nơi $g$ giảm (đỉnh núi cao) làm chu kì giảm}
\loigiai{
\begin{itemchoice}
\itemch $T \propto \sqrt{\ell}$ nên $\ell$ tăng 4 lần $\Rightarrow T$ tăng 2 lần.
\itemch Với biên độ nhỏ, $T$ không phụ thuộc biên độ.
\itemch Chu kì chỉ phụ thuộc $\ell$ và $g$, cùng $\ell$ nên cùng $T$.
\itemch $g$ giảm $\Rightarrow T$ tăng, không giảm.
\end{itemchoice}
}
\end{ex}
